% !TEX root = ../main.tex

\section{Motivation}
\label{sec:motivation}
\frame[plain]{\sectionpage}

%----------------------------------------
\begin{frame}{Synthetic financial time series: a hard problem}
    \begin{columns}[T]
    \begin{column}{.55\textwidth}
        \textbf{Why generate synthetic returns?}
        \begin{itemize}
            \item Stress-testing portfolio strategies
            \item Backtesting under counterfactual regimes
            \item Risk modelling: VaR, Expected Shortfall
            \item ML training-data augmentation
        \end{itemize}
        \medskip

        \textbf{Why is it hard?}
        \begin{itemize}
            \item Returns are non-stationary, fat-tailed,
            volatility-clustered
            \item Cross-asset dependencies are non-Gaussian
            \item Historical data is scarce relative to the
            joint distribution's dimensionality
        \end{itemize}
    \end{column}
    \hfill
    \begin{column}{.40\textwidth}
        \begin{block}{Stylised facts}
            \begin{itemize}
                \item $\sim 0$ ACF in returns
                \item Strong positive ACF in $|r_t|$ and $r_t^2$
                (volatility clustering)
                \item Excess kurtosis: heavy tails
                \item Slight negative skew
                \item Asymmetric tail dependence
            \end{itemize}
        \end{block}
    \end{column}
    \end{columns}
\end{frame}

%----------------------------------------
\begin{frame}{The quantum hypothesis}
    \begin{columns}[T]
    \begin{column}{.55\textwidth}
        \textbf{Theoretical motivation}\\
        An $n$-qubit circuit parametrises a unitary operator on a
        $\mathbf{2^n}$-dimensional Hilbert space using only
        $\mathcal{O}(\text{poly}(n))$ gates and parameters.
        \medskip

        Through \emph{entanglement}, the joint distribution of
        measurement outcomes can encode dependence structure that
        no product of independent classical operations can
        reproduce~\cite{Zoufal2019QGAN, Coyle2021QvsClassical}.
        \medskip

        \textbf{Research question}\\
        On a controlled multi-asset benchmark, does replacing the
        generator of a GAN with a small variational quantum
        circuit produce measurably different output statistics
        --- and at what parameter cost?
    \end{column}
    \hfill
    \begin{column}{.40\textwidth}
        \begin{block}{Project scope}
            \begin{itemize}
                \item Univariate \& multivariate ($\leq 5$
                assets)
                \item Vanilla GAN training~\cite{Goodfellow2014GAN}
                \item Daily log-returns, 20-day windows
                \item \smi, \dax, EuroStoxx, FTSE, \sap
                \item Simulator: PennyLane
                \cite{Bergholm2018PennyLane}
            \end{itemize}
        \end{block}
        \medskip

        Out of scope: real hardware, $> 5$ assets.
    \end{column}
    \end{columns}
\end{frame}
